\subsection{Impact of env. changes to localization ~ 50\%}
\textbf{Lifelong SLAM }

\textbf{AMCL (+ parameter tuning)}

\textbf{Pose tracking (KF and ICP)}

\subsection{Change detection ~ 50\%}
\textbf{Distance metrics, Map to map and map to sensor comparison}

Long-term_Localization_of_Mobile_Robots_in_Dynamic_Changing_Environments.pdf\newline
For each beam in a laser scan, they calculate [[Z-scores and Mahalanobis distances|Mahalanobis distance]] $d_m$, which can be thought of as the observed range difference 
(expected range vs measured) expressed as a fraction of the noise standard deviation. This is then converted to a weight w
by $w = e^{(-d)}$ (small distance -> 1, large distance -> 0). The Matching degree is then the average weight for all beams in that scan. 

The weight w is the same weight calculated by AMCL in the default Likelihood-field sensor model mode. 
The difference is that AMCL takes the product of all weights without any normalization. 
